% Options for packages loaded elsewhere
\PassOptionsToPackage{unicode}{hyperref}
\PassOptionsToPackage{hyphens}{url}
\PassOptionsToPackage{dvipsnames,svgnames,x11names}{xcolor}
%
\documentclass[
  11pt,
  a4paper,
]{article}
\usepackage{amsmath,amssymb}
\usepackage{iftex}
\ifPDFTeX
  \usepackage[T1]{fontenc}
  \usepackage[utf8]{inputenc}
  \usepackage{textcomp} % provide euro and other symbols
\else % if luatex or xetex
  \usepackage{unicode-math} % this also loads fontspec
  \defaultfontfeatures{Scale=MatchLowercase}
  \defaultfontfeatures[\rmfamily]{Ligatures=TeX,Scale=1}
\fi
\usepackage{lmodern}
\ifPDFTeX\else
  % xetex/luatex font selection
\fi
% Use upquote if available, for straight quotes in verbatim environments
\IfFileExists{upquote.sty}{\usepackage{upquote}}{}
\IfFileExists{microtype.sty}{% use microtype if available
  \usepackage[]{microtype}
  \UseMicrotypeSet[protrusion]{basicmath} % disable protrusion for tt fonts
}{}
\makeatletter
\@ifundefined{KOMAClassName}{% if non-KOMA class
  \IfFileExists{parskip.sty}{%
    \usepackage{parskip}
  }{% else
    \setlength{\parindent}{0pt}
    \setlength{\parskip}{6pt plus 2pt minus 1pt}}
}{% if KOMA class
  \KOMAoptions{parskip=half}}
\makeatother
\usepackage{xcolor}
\usepackage[margin=2.25cm]{geometry}
\usepackage{color}
\usepackage{fancyvrb}
\newcommand{\VerbBar}{|}
\newcommand{\VERB}{\Verb[commandchars=\\\{\}]}
\DefineVerbatimEnvironment{Highlighting}{Verbatim}{commandchars=\\\{\}}
% Add ',fontsize=\small' for more characters per line
\newenvironment{Shaded}{}{}
\newcommand{\AlertTok}[1]{\textcolor[rgb]{1.00,0.00,0.00}{\textbf{#1}}}
\newcommand{\AnnotationTok}[1]{\textcolor[rgb]{0.38,0.63,0.69}{\textbf{\textit{#1}}}}
\newcommand{\AttributeTok}[1]{\textcolor[rgb]{0.49,0.56,0.16}{#1}}
\newcommand{\BaseNTok}[1]{\textcolor[rgb]{0.25,0.63,0.44}{#1}}
\newcommand{\BuiltInTok}[1]{\textcolor[rgb]{0.00,0.50,0.00}{#1}}
\newcommand{\CharTok}[1]{\textcolor[rgb]{0.25,0.44,0.63}{#1}}
\newcommand{\CommentTok}[1]{\textcolor[rgb]{0.38,0.63,0.69}{\textit{#1}}}
\newcommand{\CommentVarTok}[1]{\textcolor[rgb]{0.38,0.63,0.69}{\textbf{\textit{#1}}}}
\newcommand{\ConstantTok}[1]{\textcolor[rgb]{0.53,0.00,0.00}{#1}}
\newcommand{\ControlFlowTok}[1]{\textcolor[rgb]{0.00,0.44,0.13}{\textbf{#1}}}
\newcommand{\DataTypeTok}[1]{\textcolor[rgb]{0.56,0.13,0.00}{#1}}
\newcommand{\DecValTok}[1]{\textcolor[rgb]{0.25,0.63,0.44}{#1}}
\newcommand{\DocumentationTok}[1]{\textcolor[rgb]{0.73,0.13,0.13}{\textit{#1}}}
\newcommand{\ErrorTok}[1]{\textcolor[rgb]{1.00,0.00,0.00}{\textbf{#1}}}
\newcommand{\ExtensionTok}[1]{#1}
\newcommand{\FloatTok}[1]{\textcolor[rgb]{0.25,0.63,0.44}{#1}}
\newcommand{\FunctionTok}[1]{\textcolor[rgb]{0.02,0.16,0.49}{#1}}
\newcommand{\ImportTok}[1]{\textcolor[rgb]{0.00,0.50,0.00}{\textbf{#1}}}
\newcommand{\InformationTok}[1]{\textcolor[rgb]{0.38,0.63,0.69}{\textbf{\textit{#1}}}}
\newcommand{\KeywordTok}[1]{\textcolor[rgb]{0.00,0.44,0.13}{\textbf{#1}}}
\newcommand{\NormalTok}[1]{#1}
\newcommand{\OperatorTok}[1]{\textcolor[rgb]{0.40,0.40,0.40}{#1}}
\newcommand{\OtherTok}[1]{\textcolor[rgb]{0.00,0.44,0.13}{#1}}
\newcommand{\PreprocessorTok}[1]{\textcolor[rgb]{0.74,0.48,0.00}{#1}}
\newcommand{\RegionMarkerTok}[1]{#1}
\newcommand{\SpecialCharTok}[1]{\textcolor[rgb]{0.25,0.44,0.63}{#1}}
\newcommand{\SpecialStringTok}[1]{\textcolor[rgb]{0.73,0.40,0.53}{#1}}
\newcommand{\StringTok}[1]{\textcolor[rgb]{0.25,0.44,0.63}{#1}}
\newcommand{\VariableTok}[1]{\textcolor[rgb]{0.10,0.09,0.49}{#1}}
\newcommand{\VerbatimStringTok}[1]{\textcolor[rgb]{0.25,0.44,0.63}{#1}}
\newcommand{\WarningTok}[1]{\textcolor[rgb]{0.38,0.63,0.69}{\textbf{\textit{#1}}}}
\setlength{\emergencystretch}{3em} % prevent overfull lines
\providecommand{\tightlist}{%
  \setlength{\itemsep}{0pt}\setlength{\parskip}{0pt}}
\setcounter{secnumdepth}{-\maxdimen} % remove section numbering
\ifLuaTeX
  \usepackage{selnolig}  % disable illegal ligatures
\fi
\usepackage{bookmark}
\IfFileExists{xurl.sty}{\usepackage{xurl}}{} % add URL line breaks if available
\urlstyle{same}
\hypersetup{
  colorlinks=true,
  linkcolor={blue},
  filecolor={Maroon},
  citecolor={Blue},
  urlcolor={blue},
  pdfcreator={LaTeX via pandoc}}

\author{}
\date{}

\usepackage{microtype}
\usepackage{fancyhdr}
\pagestyle{fancy}
\fancyhf{}
\lhead{Documentation-Driven Threat Analysis}
\rhead{BA4 closure}
\cfoot{\thepage}
\setlength{\headheight}{14pt}
\setlength{\emergencystretch}{4em}
\hbadness=10000
\hfuzz=20pt
\sloppy
\begin{document}

\section{DDTA BA4 projection boundary, traceability, interpretation and
coverage contract -
R1}\label{ddta-ba4-projection-boundary-traceability-interpretation-and-coverage-contract---r1}

\textbf{Status:} CLOSED BY BA4-T3 / ACCEPTED FOR CURRENT THESIS SCOPE
\textbf{R24 alignment note:} the BA4 projection contract remains retained. Statements below that BA2 reopening was not triggered and that BA5 was the next phase are closure-time dispositions, not the current R24 execution sequence. Current forward state is governed by \texttt{DDTA\_R24\_DECISION\_RULE\_CHECKPOINT.md}. This alignment does not reopen BA4.


\textbf{Closure baseline reviewed:}
\texttt{dcb4605448de4ac5331f10ff090a9f2ab677427e}

\textbf{Closed dependencies:} BA0 responsibility boundary; BA1
\texttt{BAReferent\ +\ BAProposition}; BA2 proposition semantics; BA3
provenance/derivation/lifecycle/change contract; BA4-T1
projection-boundary lower bound; BA4-T2 interpretation/coverage
refinement.

\subsection{1. Closure statement}\label{closure-statement}

BA4 closes the smallest current-scope contract by which one accepted
Base Analysis can support multiple human and methodology-oriented
projections without allowing a projection to become project authority or
consumer taxonomy to redefine shared BA meaning.

The final contract preserves six properties simultaneously:

\begin{enumerate}
\def\labelenumi{\arabic{enumi}.}
\tightlist
\item
  every meaning-bearing projection item can be traced to accepted BA
  meaning for one BA baseline;
\item
  shared semantic rendering cannot strengthen, invert or erase material
  BA meaning;
\item
  method-owned interpretation may be richer than BA rendering but
  remains explicitly downstream and review-reproducible;
\item
  projection lossiness and completeness are stated relative to an
  explicit declared scope and coverage mode;
\item
  stale, diagnostic and pending BA material may be surfaced only under
  an explicit qualification policy that preserves its
  non-current/unresolved character; and
\item
  projection rebuild and cross-projection comparison use BA
  identity/trace and BA3 continuity rather than a second project
  lifecycle or universal projection ontology.
\end{enumerate}

BA4 remains representation-independent. These are semantic
responsibilities, not mandatory classes, graph nodes, tables, JSON
records, UI objects or one executable projection language.

\subsection{2. Final projection lower
bound}\label{final-projection-lower-bound}

The final current-scope lower bound separates a revisioned projection
definition, a baseline-scoped materialization and meaning-bearing
projection items.

\begin{Shaded}
\begin{Highlighting}[]
\NormalTok{BAProjectionDescriptor                         [projection metadata; NOT BAE]}
\NormalTok{|{-} projectionKey                         1}
\NormalTok{|{-} projectionRevisionKey                 1 immutable}
\NormalTok{|{-} consumerPurpose                       1}
\NormalTok{|{-} selectionCoverageContract             1}
\NormalTok{|    |{-} eligibleBAScope                  1}
\NormalTok{|    |{-} coverageMode                     1}
\NormalTok{|    |    EXHAUSTIVE\_FOR\_DECLARED\_SCOPE}
\NormalTok{|    |    | SELECTIVE}
\NormalTok{|    \textasciigrave{}{-} qualificationPolicy              1}
\NormalTok{|{-} mappingContract                       1}
\NormalTok{\textasciigrave{}{-} interpretationRuleDescriptor          0..*}
\NormalTok{     |{-} ruleKey                          1}
\NormalTok{     |{-} inputRoleContract                1..*}
\NormalTok{     |{-} applicabilityContract            1}
\NormalTok{     |{-} outputKindContract               1}
\NormalTok{     \textasciigrave{}{-} normativeDefinition              1}

\NormalTok{BAProjectionMaterialization                    [derived artifact; NOT BAE]}
\NormalTok{|{-} projectionRef                         1 projectionKey@revision}
\NormalTok{|{-} sourceBABaselineKey                   1}
\NormalTok{\textasciigrave{}{-} item                                  0..*}

\NormalTok{BAProjectionItem                               [projection{-}local; NOT BAE]}
\NormalTok{|{-} projectionItemKey                     1 local to materialization}
\NormalTok{|{-} projectionOwnedKind                   1}
\NormalTok{|{-} traceBinding                           1..*}
\NormalTok{|    |{-} baElementRef                      1 BA element @ source BA baseline}
\NormalTok{|    \textasciigrave{}{-} traceRoleKey                      0..1}
\NormalTok{|         [required where contribution roles differ]}
\NormalTok{|{-} sharedSemanticRendering               0..1}
\NormalTok{|{-} methodOwnedInterpretation             0..1}
\NormalTok{\textasciigrave{}{-} interpretationRuleRef                 0..1}
\NormalTok{     [required for meaning{-}bearing method{-}owned interpretation]}
\end{Highlighting}
\end{Shaded}

An implementation may combine or derive these responsibilities
physically. The semantic distinctions remain mandatory where applicable.

\subsection{3. Authority and ownership
boundary}\label{authority-and-ownership-boundary}

The authority chain is fixed:

\begin{Shaded}
\begin{Highlighting}[]
\NormalTok{governed documentation}
\NormalTok{    {-}\textgreater{} accepted Base Analysis}
\NormalTok{        {-}\textgreater{} projection materialization}
\NormalTok{            {-}\textgreater{} optional method{-}owned interpretation}
\end{Highlighting}
\end{Shaded}

A projection is derived consumer state. It is never the source of a
grounded BA element and never project documentation authority.

If a projection exposes a possible gap, missing fact or corrective idea,
that output remains downstream candidate material. Project truth changes
only through governed documentation and a subsequent accepted BA
baseline.

A method-owned interpretation does not become shared BA meaning merely
because it is reproducible, useful or supported by several BA elements.

\subsection{4. Immutable projection
revision}\label{immutable-projection-revision}

A projection reference resolves to one immutable, inspectable revision.
A stable key without revision identity is insufficient because
selection, qualification, mapping and interpretation rules can change
while the BA baseline remains fixed.

Therefore:

\begin{Shaded}
\begin{Highlighting}[]
\NormalTok{same BA baseline}
\NormalTok{+ projection descriptor R1}
\NormalTok{    !=}
\NormalTok{same BA baseline}
\NormalTok{+ projection descriptor R2}
\end{Highlighting}
\end{Shaded}

The two materializations may differ without implying any change in
project meaning.

The projection descriptor is consumer-specific and is distinct from the
BA3 method-neutral derivation-rule registry. Merging those registries is
rejected because it would allow method semantics to become shared BA
derivation semantics.

\subsection{5. Coverage contract}\label{coverage-contract}

\subsubsection{5.1 Eligible BA scope}\label{eligible-ba-scope}

\texttt{eligibleBAScope} states which accepted BA meanings are
candidates for inclusion under the projection purpose. It may be
expressed through operators, semantic kinds, named rules, inspectable
queries or another bounded representation. BA4 does not require a
universal query language.

\subsubsection{5.2 Coverage mode}\label{coverage-mode}

The final contract retains two semantically distinct coverage modes:

\begin{Shaded}
\begin{Highlighting}[]
\NormalTok{EXHAUSTIVE\_FOR\_DECLARED\_SCOPE}
\NormalTok{SELECTIVE}
\end{Highlighting}
\end{Shaded}

For \texttt{EXHAUSTIVE\_FOR\_DECLARED\_SCOPE}, every BA element that is
eligible under the declared scope and qualification policy is expected
to be represented according to the mapping contract.

For \texttt{SELECTIVE}, omission of otherwise eligible BA meaning is
permitted and carries no completeness claim.

The mode must be semantically explicit. It may be serialized as an enum,
fixed textual contract or equivalent inspectable representation.

\subsubsection{5.3 Omission semantics are derived, not independently
configurable}\label{omission-semantics-are-derived-not-independently-configurable}

BA4-T3 removes \texttt{omissionSemantics} as a separate mandatory
responsibility.

Its meaning is closed by \texttt{coverageMode}:

\begin{Shaded}
\begin{Highlighting}[]
\NormalTok{EXHAUSTIVE\_FOR\_DECLARED\_SCOPE}
\NormalTok{  eligible + qualified BA element omitted}
\NormalTok{      {-}\textgreater{} projection materialization defect}

\NormalTok{SELECTIVE}
\NormalTok{  eligible BA element omitted}
\NormalTok{      {-}\textgreater{} no project{-}semantic conclusion}
\end{Highlighting}
\end{Shaded}

In neither case does omission mean that the project fact is false,
absent from BA or absent from governed documentation.

Allowing an independent omission policy would permit a projection to
reinterpret absence inconsistently with its declared completeness mode,
so the extra field is rejected as unnecessary configurability.

\subsection{6. Qualification policy}\label{qualification-policy}

Coverage and qualification are orthogonal responsibilities even when
physically stored in one coverage contract.

\texttt{eligibleBAScope} answers:

\begin{quote}
What BA meaning is relevant to this projection purpose?
\end{quote}

\texttt{coverageMode} answers:

\begin{quote}
Does the projection claim complete representation of that declared
scope?
\end{quote}

\texttt{qualificationPolicy} answers:

\begin{quote}
Which BA review/freshness/origin qualifications are admissible as
projection inputs and how must they remain visible in consumed meaning?
\end{quote}

A current-state projection may select only \texttt{ACCEPTED\ +\ CURRENT}
BA meaning. A review-oriented projection may intentionally include
\texttt{STALE}, \texttt{PENDING\_REVIEW} or
\texttt{DIAGNOSTIC\_UNRESOLVED} material.

The critical invariant is:

\begin{Shaded}
\begin{Highlighting}[]
\NormalTok{accepted diagnostic != accepted project fact}
\NormalTok{stale BA basis       != current project truth}
\end{Highlighting}
\end{Shaded}

Qualification may be resolved through the traced BA element and BA3
metadata; BA4 does not duplicate BA3 state as a second lifecycle
authority.

Collapsing qualification into rendering alone is rejected because
non-visual consumers can still misconsume stale or diagnostic basis. The
admissibility policy is therefore part of the projection contract, while
visual badges or labels remain presentation choices.

\subsection{7. Traceability lower bound}\label{traceability-lower-bound}

Every meaning-bearing projection item traces to one or more BA element
materializations belonging to \texttt{sourceBABaselineKey}.

The minimum drill-down path is:

\begin{Shaded}
\begin{Highlighting}[]
\NormalTok{projection item}
\NormalTok{    {-}\textgreater{} BAElementRef@sourceBABaseline}
\NormalTok{        {-}\textgreater{} BA3 provenance}
\NormalTok{            {-}\textgreater{} governed source}
\end{Highlighting}
\end{Shaded}

Duplicating \texttt{GovernedSourceRef} in every projection item is not
required. A projection may cache source labels or excerpts for
navigation, but BA3 remains the authoritative source-lineage layer.

A projection that bypasses accepted BA and rereads governed prose to
reconstruct already-shared meaning fails the BA4 shared projection
contract. If a consumer repeatedly needs shared methodology-neutral
meaning that is absent from BA, that is evidence for the smallest
BA1/BA2/BA3 reopen - not permission for the projection to silently
repair BA.

\subsection{8. Role-bound trace}\label{role-bound-trace}

A bare BA reference set is sufficient only when the referenced elements
contribute in an undifferentiated way.

When several BA inputs play different semantic roles in a mapping or
interpretation, trace roles are required. Example:

\begin{Shaded}
\begin{Highlighting}[]
\NormalTok{serviceUse}
\NormalTok{  {-}\textgreater{} consumeService(...)}

\NormalTok{responsibilityBoundary}
\NormalTok{  {-}\textgreater{} assignResponsibility[negative](...)}
\end{Highlighting}
\end{Shaded}

The trace roles are projection-contract scoped. They are not new BA2
proposition roles.

BA4-T3 therefore minimizes the final shape by making
\texttt{traceRoleKey} conditional rather than mandatory on every single
trace binding.

\subsection{9. Shared semantic
rendering}\label{shared-semantic-rendering}

A projection item presented as shared project meaning must be materially
supportable by its traced BA basis without strengthening, inversion or
loss of material qualification.

Allowed transformations include:

\begin{itemize}
\tightlist
\item
  selection under the coverage contract;
\item
  human-readable ordering;
\item
  presentation labels that preserve meaning;
\item
  grouping for navigation; and
\item
  aggregation whose shared semantics are no stronger than the traced BA
  basis.
\end{itemize}

Examples:

\begin{Shaded}
\begin{Highlighting}[]
\NormalTok{transfer(...) {-}\textgreater{} "delivery" with same participants                 PASS}
\NormalTok{several explicit constraints {-}\textgreater{} visibly enumerated constraint group PASS}
\NormalTok{negative responsibility {-}\textgreater{} affirmative ownership                    REJECTED}
\NormalTok{conditional proposition {-}\textgreater{} unconditional statement                  REJECTED}
\NormalTok{C/I/provenance constraints {-}\textgreater{} "the channel is secure"               REJECTED}
\NormalTok{DIAGNOSTIC\_UNRESOLVED {-}\textgreater{} established project fact                    REJECTED}
\end{Highlighting}
\end{Shaded}

The test is semantic preservation, not byte, notation or layout
identity.

\subsection{10. Method-owned
interpretation}\label{method-owned-interpretation}

A projection may add a consumer-specific interpretation beyond simple
rendering, for example:

\begin{Shaded}
\begin{Highlighting}[]
\NormalTok{externally{-}provided{-}transport{-}exposure}
\end{Highlighting}
\end{Shaded}

or a different method may derive:

\begin{Shaded}
\begin{Highlighting}[]
\NormalTok{delegated{-}assurance{-}dependency}
\end{Highlighting}
\end{Shaded}

from overlapping BA basis.

These local categories need not agree and are not BA semantic keys.

A meaning-bearing method-owned interpretation must:

\begin{enumerate}
\def\labelenumi{\arabic{enumi}.}
\tightlist
\item
  remain explicitly consumer/method owned;
\item
  trace to its BA input basis;
\item
  identify the applicable local interpretation rule; and
\item
  remain incapable of becoming shared BA/project truth without the
  governed-document feedback path.
\end{enumerate}

\subsection{11. Interpretation rule
accountability}\label{interpretation-rule-accountability}

BA4-T3 confirms that
\texttt{projectionRevisionKey\ +\ role-bound\ BA\ trace} is insufficient
when one projection revision contains multiple interpretation rules over
overlapping inputs.

Therefore \texttt{interpretationRuleRef} remains required for
meaning-bearing method-owned interpretation.

The reference resolves within the immutable projection revision to an
inspectable consumer-owned descriptor containing at least:

\begin{Shaded}
\begin{Highlighting}[]
\NormalTok{ruleKey}
\NormalTok{inputRoleContract}
\NormalTok{applicabilityContract}
\NormalTok{outputKindContract}
\NormalTok{normativeDefinition}
\end{Highlighting}
\end{Shaded}

No universal executable DSL is required.

Review reproducibility means that an independent reviewer can determine
from the same BA baseline, same role-bound trace basis, same projection
revision and same local rule whether materially equivalent local
interpretation is justified or the rule is non-applicable.

The output remains projection-local even when reproducible.

\subsection{12. Shared rendering and method interpretation remain
distinct under
aggregation}\label{shared-rendering-and-method-interpretation-remain-distinct-under-aggregation}

BA4-T3 explicitly tries to collapse shared aggregation and method
interpretation.

The collapse fails.

If an aggregate merely reorganizes or summarizes entailed BA meaning
without a stronger claim, it may remain shared semantic rendering.

If aggregation produces a new consumer conclusion, classification,
exposure, assurance judgment or method-specific abstraction, it is
method-owned interpretation and requires local rule accountability.

Example:

\begin{Shaded}
\begin{Highlighting}[]
\NormalTok{confidentiality + integrity + provenance constraints}
\NormalTok{    {-}\textgreater{} enumerated "delivery constraints" group}
\NormalTok{       SHARED RENDERING}

\NormalTok{same basis}
\NormalTok{    {-}\textgreater{} "high{-}assurance channel"}
\NormalTok{       METHOD{-}OWNED at best, if a method rule justifies it;}
\NormalTok{       never shared BA merely by aggregation}
\end{Highlighting}
\end{Shaded}

Thus storage may combine the two fields, but ownership semantics may not
be erased.

\subsection{13. Cross-projection
consistency}\label{cross-projection-consistency}

Cross-projection consistency is intentionally weaker than taxonomy
agreement.

For two projections over the same BA baseline, the final minimum
requires:

\begin{enumerate}
\def\labelenumi{\arabic{enumi}.}
\tightlist
\item
  truthful \texttt{sourceBABaselineKey};
\item
  shared renderings supportable by each item's own BA trace;
\item
  method interpretations explicitly local and rule-accountable where
  required;
\item
  satisfaction of each projection's own coverage mode and qualification
  policy; and
\item
  comparison through BA trace rather than assumed method-category
  equivalence.
\end{enumerate}

Therefore two different method labels over overlapping BA basis may both
be valid.

A material projection inconsistency exists when a projection violates BA
meaning or its own descriptor, for example by inverting polarity,
omitting an eligible item under exhaustive coverage, presenting stale
basis as current truth, promoting local interpretation to BA authority
or claiming one BA baseline while tracing silently to another.

No BA proposition saying ``the projections conflict'' is required merely
to detect such a projection-contract violation.

\subsection{14. BA trace is the shared cross-projection
anchor}\label{ba-trace-is-the-shared-cross-projection-anchor}

BA4-T3 again tries to force \texttt{SharedProjectionConcept} identity or
method-category equivalence relations.

The attempt fails for current evidence.

Same-baseline comparison can use trace overlap and role bindings:

\begin{Shaded}
\begin{Highlighting}[]
\NormalTok{Projection A item {-}\textgreater{} BA trace set}
\NormalTok{Projection B item {-}\textgreater{} BA trace set}

\NormalTok{compare shared/overlapping/disjoint BA basis}
\end{Highlighting}
\end{Shaded}

Across baselines, BA3 continuity
(\texttt{RETAIN\ \textbar{}\ REPLACE\ \textbar{}\ RETIRE}) supplies the
semantic bridge before each projection is rebuilt.

Different labels do not establish contradiction. Same labels do not
establish shared identity.

A universal projection ontology would add method taxonomy to the shared
core without a methodology-neutral project need and is rejected.

\subsection{15. Projection-local
identity}\label{projection-local-identity}

Meaning-bearing projection items need a local address within one
materialization for trace, inspection and rendering. That identity is
projection-local and is not a \texttt{BAReferent} or
\texttt{BAProposition} identity.

Pure layout primitives do not require semantic identity.

BA4 does not require cross-baseline projection-item identity or a
BA-like item lifecycle. A method/tool may maintain its own UI/history
continuity, but this remains consumer-owned unless later evidence
demonstrates a shared DDTA responsibility.

\subsection{16. Rebuild and lifecycle
boundary}\label{rebuild-and-lifecycle-boundary}

A projection materialization is scoped to one BA baseline and one
immutable projection revision.

When BA changes:

\begin{Shaded}
\begin{Highlighting}[]
\NormalTok{projection@B0 remains historical}
\NormalTok{accepted BA@B1 becomes new projection input}
\NormalTok{projection@B1 is rebuilt}
\end{Highlighting}
\end{Shaded}

When only the projection descriptor changes:

\begin{Shaded}
\begin{Highlighting}[]
\NormalTok{projection@descriptor{-}R1 remains historical}
\NormalTok{same accepted BA}
\NormalTok{    {-}\textgreater{} build projection@descriptor{-}R2}
\end{Highlighting}
\end{Shaded}

BA4 does not duplicate BA3
\texttt{RETAIN\ \textbar{}\ REPLACE\ \textbar{}\ RETIRE} into a second
project lifecycle for projection items.

A rebuilt projection follows accepted BA meaning. A retired BA transfer
cannot remain in a current exhaustive transfer projection merely because
a local projection item existed at B0.

\subsection{17. Integrated facial M1-M4
regression}\label{integrated-facial-m1-m4-regression}

\subsubsection{M1 - Ethernet to Wi-Fi}\label{m1---ethernet-to-wi-fi}

The BA transfer/service-use meaning may retain while the realization
proposition is replaced. A flow projection that excludes realization can
remain materially stable; an assurance/realization projection can
change. Different projection deltas are consistent when each follows its
own coverage contract.

\subsubsection{M2 - external to project-owned
transport}\label{m2---external-to-project-owned-transport}

Negative responsibility may be replaced by affirmative responsibility
and external service use may retire. A local
\texttt{externally-provided-transport-exposure} interpretation becomes
non-applicable on rebuild. Another method may produce a project-owned
authority-placement item. No common method category is needed.

\subsubsection{M3 - remote to local
recognition}\label{m3---remote-to-local-recognition}

Retirement of the remote transfer removes the corresponding exchange
from a current exhaustive transfer projection. Unrelated
responsibility/contract items can remain if their BA basis remains
current.

\subsubsection{M4 - representation
conflict}\label{m4---representation-conflict}

A current-state projection excludes stale/unresolved transfer meaning
under its qualification policy. A review projection may expose the stale
transfer and diagnostic conflict only with their qualification
preserved. Neither projection may invent the post-resolution project
truth before governance resolves the documentation.

All four controls pass without a projection lifecycle ontology or BA
reopen.

\subsection{18. Order/WMS and provider
regression}\label{orderwms-and-provider-regression}

Under WMS externalization, flow/integration and assurance/responsibility
projections can change in different shapes because they cover different
BA meanings. Internal mutation propositions retired by BA disappear only
from projections whose declared scope includes them; new
service/contract meanings may appear where eligible.

Provider normalization preserves the authority chain:

\begin{Shaded}
\begin{Highlighting}[]
\NormalTok{provider raw state}
\NormalTok{    {-}\textgreater{} governed mapping / BA3 derivation}
\NormalTok{        {-}\textgreater{} accepted normalized BA meaning}
\NormalTok{            {-}\textgreater{} projection / method interpretation}
\end{Highlighting}
\end{Shaded}

A method projection may not privately map raw provider vocabulary into
supposedly shared project truth.

These controls pass without \texttt{ExternalSystem},
\texttt{TrustBoundary}, \texttt{IntegrationRisk} or other method
categories becoming BA1 types.

\subsection{19. No new BAE family or BA2 operator is
forced}\label{no-new-bae-family-or-ba2-operator-is-forced}

\texttt{BAProjectionDescriptor}, \texttt{BAProjectionMaterialization},
projection items, local interpretation rules and coverage
responsibilities exist to build, review and consume views. They do not
represent independently reusable methodology-neutral project meaning or
shared project assertions.

Therefore:

\begin{Shaded}
\begin{Highlighting}[]
\NormalTok{new BAE family for Projection                     NOT FORCED}
\NormalTok{new BAE family for ProjectionItem                 NOT FORCED}
\NormalTok{method node/interaction as BA type                REJECTED}
\NormalTok{SharedProjectionConcept                           NOT FORCED}
\NormalTok{method{-}category equivalence BA relation           NOT FORCED}
\NormalTok{new BA2 operator for projection comparison        NOT FORCED}
\NormalTok{BA1 reopen                                        NOT TRIGGERED}
\NormalTok{BA2 reopen                                        NOT TRIGGERED}
\NormalTok{BA3 reopen                                        NOT TRIGGERED}
\end{Highlighting}
\end{Shaded}

\subsection{20. Final dispositions}\label{final-dispositions}

\begin{Shaded}
\begin{Highlighting}[]
\NormalTok{projection as project authority                         REJECTED}
\NormalTok{projection derived from accepted BA                     REQUIRED}
\NormalTok{immutable projection revision                           REQUIRED}
\NormalTok{eligible BA scope                                       REQUIRED}
\NormalTok{coverageMode                                            REQUIRED}
\NormalTok{  EXHAUSTIVE\_FOR\_DECLARED\_SCOPE | SELECTIVE}
\NormalTok{independent omissionSemantics field                     REJECTED AS REDUNDANT}
\NormalTok{qualificationPolicy                                     REQUIRED}
\NormalTok{meaning{-}bearing item {-}\textgreater{} BA trace                        REQUIRED}
\NormalTok{role{-}bound trace                                        REQUIRED WHERE ROLES DIFFER}
\NormalTok{mandatory duplicate governed{-}source provenance          REJECTED}
\NormalTok{shared semantic rendering                               ALLOWED WITH PRESERVATION}
\NormalTok{method{-}owned interpretation                             ALLOWED DOWNSTREAM}
\NormalTok{method interpretation {-}\textgreater{} shared BA                      REJECTED}
\NormalTok{interpretationRuleRef                                   REQUIRED WHEN METHOD MEANING{-}BEARING}
\NormalTok{universal projection DSL                                REJECTED}
\NormalTok{shared rendering + method interpretation collapse       REJECTED}
\NormalTok{universal projection ontology                           REJECTED}
\NormalTok{cross{-}projection shared identity                        NOT FORCED}
\NormalTok{BA trace + BA3 continuity as common anchor              ACCEPTED}
\NormalTok{projection item BA{-}like lifecycle                       REJECTED AS UNNECESSARY}
\NormalTok{projection rebuild after BA change                      REQUIRED}
\NormalTok{new BAE family                                          NOT FORCED}
\NormalTok{new BA2 operator                                        NOT FORCED}
\NormalTok{BA1 / BA2 / BA3 reopen                                  NOT TRIGGERED}
\NormalTok{BA4                                                     CLOSED FOR CURRENT THESIS SCOPE}
\end{Highlighting}
\end{Shaded}

\subsection{21. Reopen criteria}\label{reopen-criteria}

Reopen only the smallest BA4 responsibility if a concrete governed
corpus/projection demonstrates one of the following:

\begin{enumerate}
\def\labelenumi{\arabic{enumi}.}
\tightlist
\item
  an explicit selective/exhaustive coverage distinction cannot express a
  required completeness contract;
\item
  qualification policy cannot prevent non-current/unresolved BA meaning
  from being consumed as current shared truth;
\item
  shared semantic rendering cannot preserve required project meaning
  without new methodology-neutral BA semantics;
\item
  non-trivial method-owned interpretation cannot be independently
  reviewed using BA trace plus an inspectable local rule revision;
\item
  two projections require a shared methodology-neutral semantic identity
  that is not recoverable from BA identity/trace;
\item
  projection rebuild requires shared cross-baseline projection-item
  lifecycle semantics rather than BA3 continuity and rebuild;
\item
  source drill-down through projection -\textgreater{} BA
  -\textgreater{} BA3 provenance loses required governed-source
  accountability; or
\item
  a consumer exposes missing shared methodology-neutral project meaning
  that cannot be represented by the closed BA1-BA3 contract.
\end{enumerate}

Do not reopen BA4 merely because a method introduces a new local
category, a projection changes rendering technology, a UI wants stable
item IDs, a projection descriptor receives a new revision, or a tool
prefers a different storage model.

\subsection{22. Closure consequence}\label{closure-consequence}

BA4 is closed for current thesis scope.

The Base Analysis milestone as a whole remains open because BA5
lexical/assistance boundaries and BA6 integrated regression/holdout
closure are still required.

The next phase may start only with the smallest BA5 lexical-boundary
trial defined by the active work plan after this closure package becomes
an official Git baseline.

\end{document}
